\documentclass[12pt,a4paper]{article}
\usepackage[spanish,es-nodecimaldot]{babel}
\usepackage[utf8]{inputenc}
\usepackage[T1]{fontenc}
\usepackage{lmodern}
\usepackage[top=1.0cm,bottom=1.0cm,left=1.3cm,right=1.3cm]{geometry}
\usepackage{amsmath,amssymb}
\usepackage[table]{xcolor}
\usepackage{booktabs,array}
\usepackage{enumitem}
\usepackage[normalem]{ulem}

\setlength{\parindent}{0pt}
\setlength{\parskip}{9pt}
\pagestyle{empty}

\newcommand{\cuerpo}{\fontsize{16.4}{20}\selectfont}

\AtBeginDocument{%
  \setlength{\abovedisplayskip}{6pt}%
  \setlength{\belowdisplayskip}{6pt}%
  \setlength{\abovedisplayshortskip}{3pt}%
  \setlength{\belowdisplayshortskip}{3pt}%
  \cuerpo
}

\begin{document}

\begin{center}
{\fontsize{21}{25}\selectfont\bfseries Prueba H} {\fontsize{19}{23}\selectfont (3 o mas muestras indep.)}
\end{center}

\begin{center}
\begin{minipage}{0.87\textwidth}
\begin{center}
{\fontsize{19}{23}\selectfont\bfseries\uline{Prueba H (Kruskal Wallis)}}
\end{center}

Alternativa no paramétrica de \textbf{Anova.} Comparo que \textbf{k} medias poblacionales (de k muestras) sean significativamente iguales. \textbf{H0:} u1=u2=...=uk

1.Unimos las k muestras y las ordenamos de forma creciente.

2. A cada valor se le indica a que muestra pertenece (I, II,k) y se le asigna un rango (posición). Si dos valores son iguales, se asigna un promedio de sus rangos.

3. Sumamos los rangos de cada muestra (R1, R2, .., Rk)

4. Aplicamos la fórmula H, si \textbf{ni$>$5.} Aceptamos o rechazamos H0
\[
H = \frac{12}{n\cdot(n+1)}\cdot \sum_{i=1}^{k} \frac{\left(R_i\right)^2}{\left(n_i\right)} - 3\cdot(n+1)
\]

Distribución chi cuadrado con \textbf{k-1 grados de libertad. k} el numero de poblaciones. \textbf{n} es la suma de todos los tamaños muest
\end{minipage}
\end{center}

\begin{minipage}[t]{0.615\textwidth}
\vspace{0pt}
\textbf{Caso Práctico:} Se evalúa la eficacia de tres tratamientos preventivos anticorrosión midiendo la profundidad máxima de las cavidades (en milésimas de pulgada) en piezas de cable sometidas a un ensayo acelerado:

\begin{itemize}[leftmargin=20pt,itemsep=20pt,topsep=14pt,parsep=0pt]
  \item \textbf{Método A} ($n_1 = 6$): 77, 54, 67, 74, 71, 66
  \item \textbf{Método B} ($n_2 = 7$): 60, 41, 59, 65, 62, 64, 52
  \item \textbf{Método C} ($n_3 = 5$): 49, 52, 69, 47, 56
\end{itemize}
\end{minipage}\hfill
\begin{minipage}[t]{0.355\textwidth}
\vspace{34pt}
{\footnotesize
$H_0$ : Las poblaciones de los tres tratamientos son idénticas ($\tilde{\mu}_A = \tilde{\mu}_B = \tilde{\mu}_C$)

\vspace{4pt}
$H_1$ : Al menos uno de los tratamientos difiere en su localización central (distinta eficacia)
}
\end{minipage}

\vspace{34pt}
\begin{center}
{\fontsize{13.6}{17}\selectfont
\renewcommand{\arraystretch}{1.25}
\setlength{\tabcolsep}{6.4pt}
\begin{tabular}{l c c c c c c c c c c c c c c c c c c}
\toprule
Valor & 41 & 47 & 49 & 52 & 52 & 54 & 56 & 59 & 60 & 62 & 64 & 65 & 66 & 67 & 69 & 71 & 74 & 77 \\
Grupo & B & C & C & B & C & A & C & B & B & B & B & B & A & A & C & A & A & A \\
\midrule
Rango & 1 & 2 & 3 & \textbf{4.5} & \textbf{4.5} & 6 & 7 & 8 & 9 & 10 & 11 & 12 & 13 & 14 & 15 & 16 & 17 & 18 \\
\bottomrule
\end{tabular}
}
\end{center}

\end{document}
